\chapter{Perspectives}
\label{ch:perspectives}

% TODO: rediger. Chapitre de cloture de la partie contributions: deux directions
% de travail ouvertes, l'une prolongeant la \cref{ch:spdnet} (architectures
% profondes sur le cone SPD), l'autre ouvrant sur un objet geometrique nouveau
% dans le memoire (la phase geometrique).

\section{Introduction}
\label{sec:perspectives-introduction}

% TODO: rediger. Annoncer les deux directions et ce qui les relie: dans les deux
% cas, un objet de la \cref{pt:context} qui n'a pas encore ete exploite jusqu'au
% bout --- le transport parallele et l'holonomie d'un cote, l'attention comme
% operateur sur une variete de l'autre.

\section{Transformers sur la variété des matrices SPD}
\label{sec:perspectives-transformers}

% TODO: rediger. Prerequis: \cref{ch:context-riemann} (metrique log-euclidienne,
% variete de Stiefel, optimisation riemannienne) et \cref{ch:spdnet}.
%
% Points a couvrir:
%   - limites des architectures SPDNet convolutives: profondeur, champ
%     receptif, absence de mecanisme de mise en relation de tokens SPD;
%   - ce que devient l'attention quand les tokens sont des points d'une
%     variete: produit scalaire vs. distance riemannienne, softmax sur des
%     dissimilarites geodesiques;
%   - la moyenne ponderee de l'attention comme moyenne de Karcher ponderee
%     (\cref{sec:context-riemann-moyennes}), et son cout;
%   - alternatives log-euclidiennes ou par tangent space mapping, et ce que
%     l'on perd en invariance affine;
%   - applications visees: series temporelles d'images RSO, EEG.

\section{Phase géométrique}
\label{sec:perspectives-phase}

% TODO: rediger. Prerequis: \cref{ch:context-riemann} (connexion de Levi-Civita,
% transport parallele) et \cref{ch:context-geometry}.
%
% Points a couvrir:
%   - holonomie et phase de Berry: la phase accumulee par transport parallele
%     le long d'une boucle est une mesure de courbure, pas un artefact;
%   - lecture signal: que devient cette phase pour un sous-espace propre
%     (variete de Grassmann) ou une covariance complexe qui derive lentement;
%   - lien avec la pseudo-covariance et la non-circularite
%     (\cref{sec:context-definition-complex});
%   - pistes d'exploitation en interferometrie RSO et en detection de
%     changement, ou la phase est aujourd'hui traitee comme une nuisance.

\wip{Perspectives: contexte scientifique}{Pourquoi maintenant, pour chacune des deux directions. SPDnet convolutif a atteint un plafond (profondeur, champ réceptif, aucun mécanisme reliant des tokens SPD); la phase accumulée par transport parallèle est traitée partout comme une nuisance et jamais comme une mesure de courbure. Environ 1,5 page.}

\wip{Perspectives: objectifs du projet}{Deux ou trois objectifs NUMÉROTÉS ET FALSIFIABLES par direction. Pas « explorer l'attention sur les variétés » mais « établir si un softmax sur des dissimilarités géodésiques préserve l'invariance affine qui motive la couche SPD, et à quel coût ». Environ 1 page.}

\wip{Perspectives: méthodologie}{Les points techniques déjà listés en commentaire, promus en prose: la moyenne de l'attention EST une moyenne de Karcher pondérée, son coût, puis les repliements log-euclidien et par espace tangent et ce que chacun sacrifie en invariance. Pour la phase: holonomie sur Grassmann, lien à la pseudo-covariance et à la non-circularité. Environ 3 à 4 pages.}

\wip{Perspectives: déroulement}{Un séquencement sur 4 à 5 ans. Quel sous-problème est un stage, lequel une thèse, lequel un post-doc. ABSENT AUJOURD'HUI, et c'est la partie qui se lit comme « sait diriger une recherche ». Environ 1,5 page.}

\wip{Perspectives: moyens}{Collaborateurs et institutions nommément; instrument de financement visé (ANR JCJC, européen, DGA); données déjà accessibles (SEAPIX, FLS, séries RSO); calcul. Plus l'encadrement projeté: combien d'étudiants, à quel niveau, sur quel sous-problème. Environ 1 page.}

\wip{Perspectives: risques et repli}{Rarement écrit, et se lit comme de la maturité plutôt que comme un doute. Si l'attention de Karcher pondérée s'avère intraitable à l'échelle, la variante log-euclidienne reste un résultat négatif publiable sur ce que coûte l'invariance. Environ 0,5 page.}

\section{Bilan}
\label{sec:perspectives-bilan}

% TODO: rediger.

